\documentclass{amsart}
% For dealing with references we use the comment environment
\usepackage{verbatim}
\newenvironment{reference}{\comment}{\endcomment}
%\newenvironment{reference}{}{}
\newenvironment{slogan}{\comment}{\endcomment}
\newenvironment{history}{\comment}{\endcomment}

% For commutative diagrams we use Xy-pic
\usepackage[all]{xy}

% We use 2cell for 2-commutative diagrams.
\xyoption{2cell}
\UseAllTwocells

% We use multicol for the list of chapters between chapters
\usepackage{multicol}

% This is generally recommended for better output
\usepackage{lmodern}
\usepackage[T1]{fontenc}

% For cross-file-references

% Package for hypertext links:
\usepackage{hyperref}

% For any local file, say "hello.tex" you want to link to please

% Theorem environments.
%
\theoremstyle{plain}
\newtheorem{theorem}[subsection]{Theorem}
\newtheorem{proposition}[subsection]{Proposition}
\newtheorem{lemma}[subsection]{Lemma}

\theoremstyle{definition}
\newtheorem{definition}[subsection]{Definition}
\newtheorem{example}[subsection]{Example}
\newtheorem{exercise}[subsection]{Exercise}
\newtheorem{situation}[subsection]{Situation}

\theoremstyle{remark}
\newtheorem{remark}[subsection]{Remark}
\newtheorem{remarks}[subsection]{Remarks}

\numberwithin{equation}{subsection}

% Macros
%
\def\lim{\mathop{\mathrm{lim}}\nolimits}
\def\colim{\mathop{\mathrm{colim}}\nolimits}
\def\Spec{\mathop{\mathrm{Spec}}}
\def\Hom{\mathop{\mathrm{Hom}}\nolimits}
\def\Ext{\mathop{\mathrm{Ext}}\nolimits}
\def\SheafHom{\mathop{\mathcal{H}\!\mathit{om}}\nolimits}
\def\SheafExt{\mathop{\mathcal{E}\!\mathit{xt}}\nolimits}
\def\Sch{\mathit{Sch}}
\def\Mor{\mathop{\mathrm{Mor}}\nolimits}
\def\Ob{\mathop{\mathrm{Ob}}\nolimits}
\def\Sh{\mathop{\mathit{Sh}}\nolimits}
\def\NL{\mathop{N\!L}\nolimits}
\def\CH{\mathop{\mathrm{CH}}\nolimits}
\def\proetale{{pro\text{-}\acute{e}tale}}
\def\etale{{\acute{e}tale}}
\def\QCoh{\mathit{QCoh}}
\def\Ker{\mathop{\mathrm{Ker}}}
\def\Im{\mathop{\mathrm{Im}}}
\def\Coker{\mathop{\mathrm{Coker}}}
\def\Coim{\mathop{\mathrm{Coim}}}

% Boxtimes
%
\DeclareMathSymbol{\boxtimes}{\mathbin}{AMSa}{"02}

%
% Macros for moduli stacks/spaces
%
\def\QCohstack{\mathcal{QC}\!\mathit{oh}}
\def\Cohstack{\mathcal{C}\!\mathit{oh}}
\def\Spacesstack{\mathcal{S}\!\mathit{paces}}
\def\Quotfunctor{\mathrm{Quot}}
\def\Hilbfunctor{\mathrm{Hilb}}
\def\Curvesstack{\mathcal{C}\!\mathit{urves}}
\def\Polarizedstack{\mathcal{P}\!\mathit{olarized}}
\def\Complexesstack{\mathcal{C}\!\mathit{omplexes}}
% \Pic is the operator that assigns to X its picard group, usage \Pic(X)
% \Picardstack_{X/B} denotes the Picard stack of X over B
% \Picardfunctor_{X/B} denotes the Picard functor of X over B
\def\Pic{\mathop{\mathrm{Pic}}\nolimits}
\def\Picardstack{\mathcal{P}\!\mathit{ic}}
\def\Picardfunctor{\mathrm{Pic}}
\def\Deformationcategory{\mathcal{D}\!\mathit{ef}}

\setlength{\emergencystretch}{3em}
\begin{document}
\title[Stacks Project, Tag 02MX]{708 --- K-zero localization and cyclic-complex kernel for Serre quotients}
\maketitle
\noindent Copyright (C) 2005--2025 Johan de Jong. Stacks Project authors. Source: \href{https://stacks.math.columbia.edu/tag/02MX}{Tag 02MX}.

\noindent Editorial difficulty note (not author text). Difficulty H2: The proof lifts presentation relations through a Serre quotient and constructs an explicit two-periodic complex from a signed multiset bijection. The precise kernel characterization involves specialized algebraic K-theory bookkeeping beyond ordinary qualifications..

\noindent Editorial provenance note (not author text). History: excerpt from revision \texttt{a0444\allowbreak 6e57e\allowbreak c1fbc\allowbreak 252a8\allowbreak 71afc\allowbreak ec775\allowbreak 2fb28\allowbreak 07b14}, homology.tex, lines 2400--2614. Reference presentation was adapted; original-author context and core support blocks were retained or added. The mathematical wording is unchanged. Licensed under GNU FDL 1.2 or later; no invariant sections or cover texts. See ../licenses/COPYING. Context blocks reproduce author definitions and supporting results; they are not additional counted problems.

\begin{definition}
\label{definition-K-zero}
Let $\mathcal{A}$ be an abelian category.
We denote $K_0(\mathcal{A})$ the
{\it zeroth $K$-group of $\mathcal{A}$}.
It is the abelian group constructed as follows.
Take the free abelian group
on the objects of $\mathcal{A}$
and for every short exact sequence
$0 \to A \to B \to C \to 0$
impose the relation $[B] - [A] - [C] = 0$.
\end{definition}

\noindent
Suppose we are given an object $M$ of an abelian category $\mathcal{A}$
and a complex of the form
\begin{equation}
\label{equation-cyclic-complex}
\xymatrix{
\ldots \ar[r] &
M \ar[r]^\varphi &
M \ar[r]^\psi &
M \ar[r]^\varphi &
M \ar[r] & \ldots
}
\end{equation}
In this situation we define
$$
H^0(M, \varphi, \psi) = \Ker(\psi)/\Im(\varphi)
, \quad\text{and}\quad
H^1(M, \varphi, \psi) = \Ker(\varphi)/\Im(\psi).
$$

\begin{definition}
\label{definition-serre-subcategory}
\begin{reference}
\href{https://stacks.math.columbia.edu/bibliography}{Bibliography: Serre\_homotopie\_classes (Condition (I) on page 259)}
\end{reference}
Let $\mathcal{A}$ be an abelian category.
\begin{enumerate}
\item A {\it Serre subcategory} of $\mathcal{A}$ is a
nonempty full subcategory $\mathcal{C}$ of $\mathcal{A}$
such that given an exact sequence\footnote{By
Definition \ref{definition-exact} this means $\Im(A \to B) = \Ker(B \to C)$.}
$$
A \to B \to C
$$
with $A, C \in \Ob(\mathcal{C})$, then also
$B \in \Ob(\mathcal{C})$.
\item A {\it weak Serre subcategory} of $\mathcal{A}$ is a nonempty
full subcategory $\mathcal{C}$ of $\mathcal{A}$ such that given an
exact sequence
$$
A_0 \to A_1 \to A_2 \to A_3 \to A_4
$$
with $A_0, A_1, A_3, A_4$ in $\mathcal{C}$, then also $A_2$ in $\mathcal{C}$.
\end{enumerate}
\end{definition}

\begin{definition}
\label{definition-exact}
Let $\mathcal{A}$ be an additive category. Consider a sequence of morphisms
$$
\ldots \to x \to y \to z \to \ldots
\quad\text{or}\quad
x_1 \to x_2 \to \ldots \to x_n
$$
in $\mathcal{A}$. We say such a sequence is a {\it complex} if the
composition of any two consecutive (drawn) arrows is zero.
If $\mathcal{A}$ is abelian then we say a complex of the first
type above is {\it exact at $y$} if $\Im(x \to y) = \Ker(y \to z)$
and we say a complex of the second kind is {\it exact at $x_i$}
where $1 < i < n$ if
$\Im(x_{i - 1} \to x_i) = \Ker(x_i \to x_{i + 1})$. We a
sequence as above is {\it exact} or is an {\it exact sequence} or is an
{\it exact complex} if it is a complex and exact at every object (in
the first case) or exact at $x_i$ for all $1 < i < n$ (in the second case).
There are variants of these notions for sequences of the form
$$
\ldots \to x_{-3} \to x_{-2} \to x_{-1}
\quad\text{and}\quad
x_1 \to x_2 \to x_3 \to \ldots
$$
A {\it short exact sequence} is an exact complex of the form
$$
0 \to A  \to B \to C \to 0.
$$
\end{definition}

\begin{definition}
\label{categories-definition-multiplicative-system}
Let $\mathcal{C}$ be a category. A set of arrows $S$ of $\mathcal{C}$ is
called a {\it left multiplicative system} if it has the following properties:
\begin{enumerate}
\item[LMS1] The identity of every object of $\mathcal{C}$ is in $S$ and
the composition of two composable elements of $S$ is in $S$.
\item[LMS2] Every solid diagram
$$
\xymatrix{
X \ar[d]_t \ar[r]_g & Y \ar@{..>}[d]^s \\
Z \ar@{..>}[r]^f & W
}
$$
with $t \in S$ can be completed to a commutative dotted square with
$s \in S$.
\item[LMS3] For every pair of morphisms $f, g : X \to Y$ and
$t \in S$ with target $X$ such that $f \circ t = g \circ t$
there exists an $s \in S$ with source $Y$ such that
$s \circ f = s \circ g$.
\end{enumerate}
A set of arrows $S$ of $\mathcal{C}$ is
called a {\it right multiplicative system}
if it has the following properties:
\begin{enumerate}
\item[RMS1] The identity of every object of $\mathcal{C}$ is in $S$ and
the composition of two composable elements of $S$ is in $S$.
\item[RMS2] Every solid diagram
$$
\xymatrix{
X \ar@{..>}[d]_t \ar@{..>}[r]_g & Y \ar[d]^s \\
Z \ar[r]^f & W
}
$$
with $s \in S$ can be completed to a commutative dotted square with
$t \in S$.
\item[RMS3] For every pair of morphisms $f, g : X \to Y$ and
$s \in S$ with source $Y$ such that $s \circ f = s \circ g$
there exists a $t \in S$ with target $X$ such that
$f \circ t = g \circ t$.
\end{enumerate}
A set of arrows $S$ of $\mathcal{C}$ is called a {\it multiplicative system}
if it is both a left multiplicative system and a right multiplicative system.
In other words, this means that MS1, MS2, MS3 hold, where
MS1 $=$ LMS1 $+$ RMS1, MS2 $=$ LMS2 $+$ RMS2, and
MS3 $=$ LMS3 $+$ RMS3. (That said, of course LMS1 $=$ RMS1
$=$ MS1.)
\end{definition}

\begin{definition}
\label{categories-definition-saturated-multiplicative-system}
Let $\mathcal{C}$ be a category and let $S$ be a multiplicative system.
We say $S$ is {\it saturated} if, in addition to MS1, MS2, MS3, we
also have
\begin{enumerate}
\item[MS4] Given three composable morphisms $f, g, h$, if
$fg, gh \in S$, then $g \in S$.
\end{enumerate}
\end{definition}

\begin{lemma}
\label{lemma-serre-subcategory-is-kernel}
Let $\mathcal{A}$ be an abelian category.
Let $\mathcal{C} \subset \mathcal{A}$ be a Serre subcategory.
There exists an abelian category $\mathcal{A}/\mathcal{C}$
and an exact functor
$$
F : \mathcal{A} \longrightarrow \mathcal{A}/\mathcal{C}
$$
which is essentially surjective and whose kernel is $\mathcal{C}$.
The category $\mathcal{A}/\mathcal{C}$ and the functor $F$ are
characterized by the following universal property: For any exact
functor $G : \mathcal{A} \to \mathcal{B}$ such that
$\mathcal{C} \subset \Ker(G)$ there exists a factorization
$G = H \circ F$ for a unique exact functor
$H : \mathcal{A}/\mathcal{C} \to \mathcal{B}$.
\end{lemma}

\begin{proof}
Consider the set of arrows of $\mathcal{A}$ defined by
the following formula
$$
S = \{f \in \text{Arrows}(\mathcal{A}) \mid
\Ker(f), \Coker(f) \in \Ob(\mathcal{C}) \}.
$$
We claim that $S$ is a multiplicative system. To prove this we have
to check MS1, MS2, MS3, see
Categories, Definition \ref{categories-definition-multiplicative-system}.

\medskip\noindent
It is clear that identities are elements of $S$. Suppose that
$f : A \to B$ and $g : B \to C$ are elements of $S$.
There are exact sequences
$$
\begin{matrix}
0 \to \Ker(f) \to \Ker(gf) \to \Ker(g) \\
\Coker(f) \to \Coker(gf) \to \Coker(g) \to 0
\end{matrix}
$$
Hence it follows that $gf \in S$. This proves MS1. (In fact, a similar
argument will show that $S$ is a saturated multiplicative system, see
Categories, Definition
\ref{categories-definition-saturated-multiplicative-system}.)

\medskip\noindent
Consider a solid diagram
$$
\xymatrix{
A \ar[d]_t \ar[r]_g & B \ar@{..>}[d]^s \\
C \ar@{..>}[r]^f & C \amalg_A B
}
$$
with $t \in S$. Set
$W = C \amalg_A B =  \Coker((t, -g) : A \to C \oplus B)$.
Then $\Ker(t) \to \Ker(s)$ is surjective and
$\Coker(t) \to \Coker(s)$ is an isomorphism. Hence
$s$ is an element of $S$. This proves LMS2 and the proof of RMS2 is dual.

\medskip\noindent
Finally, consider morphisms $f, g : B \to C$ and a morphism $s : A \to B$
in $S$ such that $f \circ s = g \circ s$. This means that
$(f - g) \circ s = 0$. In turn this means that
$I = \Im(f - g) \subset C$ is a quotient of $\Coker(s)$
hence an object of $\mathcal{C}$. Thus $t : C \to C' = C/I$ is an
element of $S$ such that $t \circ (f - g) = 0$, i.e., such that
$t \circ f = t \circ g$. This proves LMS3 and the proof of
RMS3 is dual.

\medskip\noindent
Having proved that $S$ is a multiplicative system we set
$\mathcal{A}/\mathcal{C} = S^{-1}\mathcal{A}$, and we set
$F$ equal to the localization functor $Q$. By
Lemma \ref{lemma-localization-abelian}
the category $\mathcal{A}/\mathcal{C}$ is abelian and $F$ is exact.
If $X$ is in the kernel of $F = Q$, then by
Lemma \href{https://stacks.math.columbia.edu/tag/05QF}{lemma kernel localization (Tag 05QF)}
we see that $0 : X \to Z$ is an element of $S$ and hence
$X$ is an object of $\mathcal{C}$, i.e., the kernel of
$F$ is $\mathcal{C}$.
Finally, if $G$ is as in the statement of the lemma, then $G$ turns
every element of $S$ into an isomorphism. Hence we obtain the
functor $H : \mathcal{A}/\mathcal{C} \to \mathcal{B}$ from
the universal property of localization, see
Categories, Lemma \href{https://stacks.math.columbia.edu/tag/04VG}{lemma properties left localization (Tag 04VG)}.
We still have to show the functor $H$ is exact.
To do this it suffices to show that $H$ commutes
with taking kernels and cokernels, see Lemma \href{https://stacks.math.columbia.edu/tag/010N}{lemma exact functor (Tag 010N)}.
Let $A \to B$ be a morphism in $\mathcal{A}/\mathcal{C}$.
We may represent $A \to B$ as $fs^{-1}$ where $s : A' \to A$
is in $S$ and $f : A' \to B$ an arbitrary morphism of $\mathcal{A}$.
Since $F = Q$ maps $s$ to an isomorphism in the quotient category
$\mathcal{A}/\mathcal{C}$, it suffices to show that $H$ commutes with taking
kernels and cokernels of morphisms $f : A \to B$ of $\mathcal{A}$.
But here we have $H(f) = G(f)$ and the result follows
from the fact that $G$ is exact.
\end{proof}


\begin{lemma}
\label{lemma-localization-abelian}
Let $\mathcal{A}$ be an abelian category.
\begin{enumerate}
\item If $S$ is a left multiplicative system, then
the category $S^{-1}\mathcal{A}$ has cokernels and the functor
$Q : \mathcal{A} \to S^{-1}\mathcal{A}$ commutes with them.
\item If $S$ is a right multiplicative system, then
the category $S^{-1}\mathcal{A}$ has kernels and the functor
$Q : \mathcal{A} \to S^{-1}\mathcal{A}$ commutes with them.
\item If $S$ is a multiplicative system, then the category
$S^{-1}\mathcal{A}$ is abelian and the functor
$Q : \mathcal{A} \to S^{-1}\mathcal{A}$ is exact.
\end{enumerate}
\end{lemma}

\begin{proof}
Assume $S$ is a left multiplicative system. Let $a : X \to Y$ be a morphism
of $S^{-1}\mathcal{A}$. Then $a = s^{-1}f$ for some $s : Y \to Y'$
in $S$ and $f : X \to Y'$. Since $Q(s)$ is an isomorphism we see that
the existence of $\Coker(a : X \to Y)$ is equivalent to the existence
of $\Coker(Q(f) : X \to Y')$. Since $\Coker(Q(f))$ is the
coequalizer of $0$ and $Q(f)$ we see that $\Coker(Q(f))$ is
represented by $Q(\Coker(f))$ by
Categories, Lemma \href{https://stacks.math.columbia.edu/tag/05Q2}{lemma left localization limits (Tag 05Q2)}.
This proves (1).

\medskip\noindent
Part (2) is dual to part (1).

\medskip\noindent
If $S$ is a multiplicative system, then $S$ is both a left and a right
multiplicative system. Thus we see that $S^{-1}\mathcal{A}$ has
kernels and cokernels and $Q$ commutes with kernels and cokernels.
To finish the proof of (3) we have to show that $\Coim = \Im$ in
$S^{-1}\mathcal{A}$. Again using that any arrow in $S^{-1}\mathcal{A}$
is isomorphic to an arrow $Q(f)$ we see that the result follows
from the result for $\mathcal{A}$.
\end{proof}


\begin{lemma}
\label{lemma-serre-subcategory-K-groups}
Let $\mathcal{A}$ be an abelian category.
Let $\mathcal{C} \subset \mathcal{A}$ be a Serre subcategory and
set $\mathcal{B} = \mathcal{A}/\mathcal{C}$.
\begin{enumerate}
\item The exact functors $\mathcal{C} \to \mathcal{A}$ and
$\mathcal{A} \to \mathcal{B}$ induce an exact sequence
$$
K_0(\mathcal{C}) \to
K_0(\mathcal{A}) \to
K_0(\mathcal{B}) \to
0
$$
of $K$-groups, and
\item the kernel of $K_0(\mathcal{C}) \to K_0(\mathcal{A})$ is equal
to the collection of elements of the form
$$
[H^0(M, \varphi, \psi)] - [H^1(M, \varphi, \psi)]
$$
where $(M, \varphi, \psi)$ is a complex as in (\ref{equation-cyclic-complex})
with the property that it becomes exact in $\mathcal{B}$; in other words
that $H^0(M, \varphi, \psi)$ and $H^1(M, \varphi, \psi)$ are
objects of $\mathcal{C}$.
\end{enumerate}
\end{lemma}

\begin{proof}
Proof of (1). It is clear that $K_0(\mathcal{A}) \to K_0(\mathcal{B})$
is surjective and that the composition $K_0(\mathcal{C}) \to
K_0(\mathcal{A}) \to K_0(\mathcal{B})$ is zero. Let $x \in K_0(\mathcal{A})$
be an element mapping to zero in $K_0(\mathcal{B})$. We can write
$x = [A] - [A']$ with $A, A'$ in $\mathcal{A}$ (fun exercise).
Denote $B, B'$ the corresponding objects of $\mathcal{B}$. The fact that
$x$ maps to zero in $K_0(\mathcal{B})$
means that there exists a finite set $I = I^+ \amalg I^{-}$,
for each $i \in I$ a short exact sequence
$$
0 \to B_i \to B'_i \to B''_i \to 0
$$
in $\mathcal{B}$ such that we have
$$
[B] - [B'] = \sum\nolimits_{i \in I^{+}} ([B'_i] - [B_i] - [B''_i])
-
\sum\nolimits_{i \in I^{-}} ([B'_i] - [B_i] - [B''_i])
$$
in the free abelian group on isomorphism classes of objects of $\mathcal{B}$.
We can rewrite this as
$$
[B]
+ \sum\nolimits_{i \in I^{+}} ([B_i] + [B''_i])
+ \sum\nolimits_{i \in I^{-}} [B'_i]
=
[B']
+ \sum\nolimits_{i \in I^{-}} ([B_i] + [B''_i])
+ \sum\nolimits_{i \in I^{+}} [B'_i].
$$
Since the right and left hand side should contain the same isomorphism classes
of objects of $\mathcal{B}$ counted with multiplicity, this means there should
be a bijection
$$
\tau :
\{B\} \amalg \{B_i, B''_i; i \in I^+\} \amalg \{B'_i; i \in I^-\}
\longrightarrow
\{B'\} \amalg \{B_i, B''_i; i \in I^-\} \amalg \{B'_i; i \in I^+\}
$$
such that $N$ and $\tau(N)$ are isomorphic in $\mathcal{B}$.
The proof of Lemmas \ref{lemma-serre-subcategory-is-kernel} and
\ref{lemma-localization-abelian} show that we choose for $i \in I$
a short exact sequence
$$
0 \to A_i \to A'_i \to A''_i \to 0
$$
in $\mathcal{A}$ such that $B_i, B'_i, B''_i$ are isomorphic to the
images of $A_i, A'_i, A''_i$ in $\mathcal{B}$. This implies that the
corresponding bijection
$$
\tau :
\{A\} \amalg \{A_i, A''_i; i \in I^+\} \amalg \{A'_i; i \in I^-\}
\longrightarrow
\{A'\} \amalg \{A_i, A''_i; i \in I^-\} \amalg \{A'_i; i \in I^+\}
$$
satisfies the property that $M$ and $\tau(M)$ are objects of $\mathcal{A}$
which become isomorphic in $\mathcal{B}$. This means $[M] - [\tau(M)]$
is in the image of $K_0(\mathcal{C}) \to K_0(\mathcal{A})$. Namely,
the isomorphism in $\mathcal{B}$ is given by a diagram
$M \leftarrow M' \rightarrow \tau(M)$ in $\mathcal{A}$ where both
$M' \to M$ and $M' \to \tau(M)$ have kernel and cokernel in $\mathcal{C}$.
Working backwards we conclude that $x = [A] - [A']$ is in the image
of $K_0(\mathcal{C}) \to K_0(\mathcal{A})$ and the proof of part (1)
is complete.

\medskip\noindent
Proof of (2). The proof is similar to the proof of (1) but slightly
more bookkeeping is involved. First we remark that any class of the type
$[H^0(M, \varphi, \psi)] - [H^1(M, \varphi, \psi)]$ is zero
in $K_0(\mathcal{A})$ by the following calculation
\begin{align*}
0 & = [M] - [M] \\
& =  [\Ker(\varphi)] + [\Im(\varphi)]
- [\Ker(\psi)] - [\Im(\psi)] \\
& =
[\Ker(\varphi)/\Im(\psi)] -
[\Ker(\psi)/\Im(\varphi)] \\
& = [H^1(M, \varphi, \psi)] - [H^0(M, \varphi, \psi)]
\end{align*}
as desired. Hence it suffices to show that any element in the kernel
of $K_0(\mathcal{C}) \to K_0(\mathcal{A})$ is of this form.

\medskip\noindent
Any element $x$ in $K_0(\mathcal{C})$ can be represented as the
difference $x = [P] - [Q]$ of two objects of $\mathcal{C}$ (fun exercise).
Suppose that this element maps to zero in $K_0(\mathcal{A})$.
This means that there exist
\begin{enumerate}
\item a finite set $I = I^{+} \amalg I^{-}$,
\item for $i \in I$ a short exact sequence $0 \to A_i \to B_i \to C_i \to 0$
in $\mathcal{A}$
\end{enumerate}
such that
$$
[P] - [Q] =
\sum\nolimits_{i \in I^{+}} ([B_i] - [A_i] - [C_i])
-
\sum\nolimits_{i \in I^{-}} ([B_i] - [A_i] - [C_i])
$$
in the free abelian group on the objects of $\mathcal{A}$.
We can rewrite this as
$$
[P]
+ \sum\nolimits_{i \in I^{+}} ([A_i] + [C_i])
+ \sum\nolimits_{i \in I^{-}} [B_i]
=
[Q]
+ \sum\nolimits_{i \in I^{-}} ([A_i] + [C_i])
+ \sum\nolimits_{i \in I^{+}} [B_i].
$$
Since the right and left hand side should contain the same objects
of $\mathcal{A}$ counted with multiplicity, this means there should be
a bijection $\tau$ between the terms which occur above. Set
$$
T^{+} =
\{p\}\ \amalg\ \{a, c\} \times I^{+}\ \amalg\ \{b\} \times I^{-}
$$
and
$$
T^{-} =
\{q\}\ \amalg\ \{a, c\} \times I^{-}\ \amalg\ \{b\} \times I^{+}.
$$
Set $T = T^{+} \amalg T^{-} = \{p, q\} \amalg \{a, b, c\} \times I$.
For $t \in T$ define
$$
O(t)
=
\left\{
\begin{matrix}
P & \text{if} & t = p \\
Q & \text{if} & t = q \\
A_i & \text{if} & t = (a, i) \\
B_i & \text{if} & t = (b, i) \\
C_i & \text{if} & t = (c, i)
\end{matrix}
\right.
$$
Hence we can view $\tau : T^{+} \to T^{-}$ as a bijection
such that $O(t) = O(\tau(t))$ for all $t \in T^{+}$.
Let $t^{-}_0 = \tau(p)$ and let $t^{+}_0 \in T^{+}$ be the
unique element such that $\tau(t^{+}_0) = q$.
Consider the object
$$
M^{+} = \bigoplus\nolimits_{t \in T^{+}} O(t)
$$
By using $\tau$ we see that it is equal to the object
$$
M^{-} = \bigoplus\nolimits_{t \in T^{-}} O(t)
$$
Consider the map
$$
\varphi : M^{+} \longrightarrow M^{-}
$$
which on the summand $O(t) = A_i$ corresponding to $t = (a, i)$, $i \in I^{+}$
uses the map $A_i \to B_i$ into the summand $O((b, i)) = B_i$ of $M^{-}$
and on the summand $O(t) = B_i$ corresponding to $(b, i)$, $i \in I^{-}$
uses the map $B_i \to C_i$ into the summand $O((c, i)) = C_i$ of $M^{-}$.
The map is zero on the summands corresponding to $p$
and $(c, i)$, $i \in I^{+}$.
Similarly, consider the map
$$
\psi : M^{-} \longrightarrow M^{+}
$$
which on the summand $O(t) = A_i$ corresponding to $t = (a, i)$, $i \in I^{-}$
uses the map $A_i \to B_i$ into the summand $O((b, i)) = B_i$ of $M^{+}$
and on the summand $O(t) = B_i$ corresponding to $(b, i)$, $i \in I^{+}$
uses the map $B_i \to C_i$ into the summand $O((c, i)) = C_i$ of $M^{+}$.
The map is zero on the summands corresponding to $q$ and
$(c, i)$, $i \in I^{-}$.

\medskip\noindent
Note that the kernel of $\varphi$ is equal to the direct sum of the
summand $P$ and the summands $O((c, i)) = C_i$, $i \in I^{+}$ and
the subobjects $A_i$ inside the summands $O((b, i)) = B_i$, $i \in I^{-}$.
The image of $\psi$ is equal to the direct sum of the
summands $O((c, i)) = C_i$, $i \in I^{+}$ and
the subobjects $A_i$ inside the summands $O((b, i)) = B_i$, $i \in I^{-}$.
In other words we see that
$$
P \cong \Ker(\varphi)/\Im(\psi).
$$
In exactly the same way we see that
$$
Q \cong \Ker(\psi)/\Im(\varphi).
$$
Since as we remarked above the existence of the bijection
$\tau$ shows that $M^{+} = M^{-}$ we see that the lemma follows.
\end{proof}
\end{document}
